\documentclass[a4paper,12pt]{report}

% Pakiety językowe i czcionki
\usepackage[utf8]{inputenc}
\usepackage[T1]{fontenc}
\usepackage[polish]{babel}

% Pakiety formatujące i graficzne
\usepackage{geometry}
\geometry{a4paper, left=2.5cm, right=2.5cm, top=2.5cm, bottom=2.5cm}
\usepackage{hyperref}
\usepackage{tabularx}
\usepackage{booktabs}
\usepackage{xcolor}
\usepackage{listings}
\usepackage{titlesec}
\usepackage{makeidx}

% Pakiety do diagramów
\usepackage{tikz}
\usetikzlibrary{positioning, shapes.multipart, arrows.meta, backgrounds}

% Ustawienia linków i kolorów
\hypersetup{
    colorlinks=true,
    linkcolor=blue,
    filecolor=magenta,      
    urlcolor=cyan,
    pdftitle={Dokumentacja Techniczna Systemu Medycznego},
}

% Definicje kolorów dla statusów i ról
\definecolor{statusPlanned}{HTML}{1976D2}
\definecolor{statusCompleted}{HTML}{2E7D32}
\definecolor{statusCancelled}{HTML}{D32F2F}

\makeindex

\title{
    \vspace{2cm}
    \Huge \textbf{Dokumentacja Techniczna Systemu Zarządzania Kliniką}\\
    \vspace{1cm}
    \Large Architektura, API oraz Wymagania Funkcjonalne
    \vspace{2cm}
}
\author{Jakub Peska}
\date{\today}

\begin{document}

\maketitle

\tableofcontents
\newpage

% ---------------------------------------------------------
% ROZDZIAŁ 1
% ---------------------------------------------------------
\chapter{Wprowadzenie}
\section{Cel dokumentu}
Niniejszy dokument stanowi autorytatywną specyfikację techniczną systemu do zarządzania placówką medyczną. Obejmuje on zbiór wymagań funkcjonalnych, model relacyjny bazy danych oraz interfejs programistyczny aplikacji (API) dla backendu. 

\section{Założenia architektoniczne}
Fundamentem architektury opisywanego systemu jest absolutna transparentność wszystkich procesów przetwarzania danych. System został zaprojektowany z myślą o zapewnieniu pełnej widoczności realizowanych operacji, co niweluje zjawisko "ukrytych błędów" lub nieautoryzowanych modyfikacji. Narzędziem realizującym to założenie jest wbudowany, kuloodporny \index{AuditService} \texttt{AuditService}, który działa analogicznie do czarnej skrzynki, utrwalając każdy stan encji przed i po jej modyfikacji.

% ---------------------------------------------------------
% ROZDZIAŁ 2
% ---------------------------------------------------------
\chapter{Wymagania Funkcjonalne}
\section{Wstęp}
Moduł ten opisuje wymagania funkcjonalne postawione przed aplikacją ze szczególnym uwzględnieniem zarządzania zasobami ludzkimi (pacjenci, lekarze) oraz wizytami, a także definiuje zasady bezpieczeństwa (autoryzacja na podstawie ról).

\section{Przypadki Użycia i Zachowanie Systemu}

\subsection{Zarządzanie Wizytami (Appointments/Visits)}
\begin{itemize}
    \item \textbf{Modyfikacja Wizyty:} Uprawnieni użytkownicy (Rola 0 - Administrator, Rola 2 - Recepcja) mogą edytować istniejące wizyty za pomocą dedykowanych okien dialogowych (\index{Material Design} Angular Material).
    \item \textbf{Zarządzanie statusem:} Wizyty posiadają zdefiniowaną maszynę stanów: \texttt{PLANNED}, \texttt{COMPLETED}, \texttt{CANCELLED}. Na warstwie prezentacji stany te są wyróżniane za pomocą dedykowanej palety barw.
    \item \textbf{Śledzenie zmian:} Każda edycja wizyty jest przechwytywana, a stan poprzedni i nowy jest utrwalany w logach audytu.
\end{itemize}

\subsection{Zarządzanie Lekarzami}
\begin{itemize}
    \item \textbf{Podgląd harmonogramu:} Użytkownik posiada dostęp do kalendarza pracy wybranego lekarza.
    \item \textbf{Usuwanie konta lekarza:} Operacja usunięcia lekarza z systemu wymusza zastosowanie kaskadowego usuwania przypisanych mu wizyt. Zapobiega to naruszeniu więzów spójności (\index{Foreign Key} Foreign Key Constraint). Operacja usunięcia całego węzła jest raportowana w audycie. Wyłącznie administratorzy (Rola 0) posiadają uprawnienia do tej akcji.
\end{itemize}

\subsection{Zarządzanie Pacjentami}
\begin{itemize}
    \item \textbf{Zarządzanie danymi:} Edycja profilu pacjenta (imię, nazwisko, PESEL, numer telefonu).
    \item \textbf{Zabezpieczenie przed duplikacją:} PESEL posiada unikalny klucz (\texttt{UNIQUE}), wymuszający weryfikację na poziomie silnika bazy danych.
\end{itemize}

\subsection{System Audytu (Blackbox) i Inicjalizacja}
\begin{itemize}
    \item \textbf{Logowanie Akcji:} Moduł bezwzględnie rejestrujący żądania modyfikujące stan (PUT, DELETE). Zapisuje kto, kiedy, dlaczego oraz co zmienił.
    \item \textbf{Masowe zasilanie danymi (Seeder):} System jest zoptymalizowany pod obciążenia. Moduł \texttt{createdb.mjs} integruje \index{Faker} \texttt{FakerPL} w celu transakcyjnego zasilenia bazy minimum 10 000 pacjentów oraz 50 000 wizyt, generując potężny wolumen danych w kilka sekund dzięki instrukcjom \texttt{BEGIN} i \texttt{COMMIT}.
\end{itemize}

\section{Podsumowanie modułu}
System realizuje pełen cykl życia obiektów medycznych, kładąc szczególny nacisk na asynchroniczne odświeżanie interfejsu za pomocą mechanizmów detekcji zmian (ChangeDetection) w frameworku Angular oraz na spójność danych.

% ---------------------------------------------------------
% ROZDZIAŁ 3
% ---------------------------------------------------------
\chapter{Diagram Bazy Danych}
\section{Wstęp}
Silnikiem bazodanowym systemu jest relacyjna baza \index{SQLite} \texttt{SQLite}. Skrypt inicjalizacyjny zapewnia budowę struktur opartych o rygorystyczne więzy integralności (Foreign Keys = ON). 

\section{Schemat ERD (Entity-Relationship Diagram)}

\begin{figure}[htbp]
\centering
\begin{tikzpicture}[
    node distance=2cm and 2cm,
    table/.style={
        rectangle split,
        rectangle split parts=2,
        draw,
        rounded corners=2pt,
        align=left,
        text width=4.5cm,
        font=\sffamily\footnotesize
    },
    header/.style={
        fill=blue!10,
        font=\sffamily\bfseries\small
    }
]

% Tabela: Patients
\node[table] (patients) {
    \nodepart{one} \textbf{patients}
    \nodepart{two}
    \textbf{id}: INTEGER (PK) \\
    firstname: TEXT \\
    lastname: TEXT \\
    pesel: TEXT (UNIQUE) \\
    phone: TEXT
};

% Tabela: Doctors
\node[table, right=2cm of patients] (doctors) {
    \nodepart{one} \textbf{doctors}
    \nodepart{two}
    \textbf{id}: INTEGER (PK) \\
    firstname: TEXT \\
    lastname: TEXT \\
    specialization: TEXT
};

% Tabela: Visits
\node[table, below=2.5cm of patients, xshift=3.5cm] (visits) {
    \nodepart{one} \textbf{visits}
    \nodepart{two}
    \textbf{id}: INTEGER (PK) \\
    patientId: INTEGER (FK) \\
    doctorId: INTEGER (FK) \\
    visitDate: TEXT \\
    room: TEXT \\
    status: TEXT (DEFAULT 'PLANNED')
};

% Tabela: Audit Logs
\node[table, below=2cm of visits] (audit) {
    \nodepart{one} \textbf{audit\_logs}
    \nodepart{two}
    \textbf{id}: INTEGER (PK) \\
    user\_id: INTEGER \\
    user\_email: TEXT \\
    action: TEXT \\
    entity\_type: TEXT \\
    entity\_id: INTEGER \\
    old\_data: TEXT \\
    new\_data: TEXT \\
    timestamp: DATETIME \\
    reason: TEXT
};

% Relacje
\draw[-{Latex[length=3mm]}] (visits.north) -- +(0,0.5) -| (patients.south) node[pos=0.25, left, font=\scriptsize] {1..N};
\draw[-{Latex[length=3mm]}] (visits.north) -- +(0,0.5) -| (doctors.south) node[pos=0.25, right, font=\scriptsize] {1..N};

\draw[dashed, gray] (visits.east) -| +(1.5,-2) |- (audit.east) node[pos=0.25, right, font=\scriptsize, align=center] {Rejestr \\ zmian (log)};

\end{tikzpicture}
\caption{Struktura relacyjna bazy danych SQLite}
\label{fig:erd}
\end{figure}

\section{Szczegóły struktury i indeksowanie}
Tabela \texttt{audit\_logs} posiada indeks \texttt{idx\_audit\_entity} obejmujący kolumny \texttt{entity\_type} oraz \texttt{entity\_id}. Ograniczenie to gwarantuje subsekundowe czasy odpowiedzi podczas generowania historii zmian dla danego zasobu (pacjenta lub wizyty), nawet przy dziesiątkach tysięcy wpisów. Zastosowano konwencję nazewnictwa kluczy obcych camelCase (np. \texttt{patientId}) dla ujednolicenia z interfejsem warstwy backendu (Node.js/Express).

% ---------------------------------------------------------
% ROZDZIAŁ 4
% ---------------------------------------------------------
\chapter{Specyfikacja API Backendu}
\section{Wstęp}
Warstwa serwerowa została zrealizowana w środowisku Node.js z wykorzystaniem frameworka Express. Komunikacja odbywa się w standardzie \index{REST API} REST, a ładunki użyteczne wymieniane są w formacie JSON. Wszystkie ścieżki podlegają uwierzytelnieniu za pośrednictwem odpowiednich middleware (funkcja \texttt{requireAuth}).

\section{Kluczowe Endpointy}

\subsection{Aktualizacja Wizyty}
Służy do edycji parametrów zaplanowanej lub zakończonej wizyty. \index{Endpoint} Endpoint wbudowuje mechanizm automatycznego zapisu do audytu.\\

\noindent
\begin{tabularx}{\textwidth}{|l|X|}
\hline
\textbf{Metoda HTTP} & \texttt{PUT} \\ \hline
\textbf{Ścieżka} & \texttt{/api/visits/:id} \\ \hline
\textbf{Wymagane role} & \texttt{0} (Admin), \texttt{2} (Recepcja) \\ \hline
\textbf{Parametry ścieżki} & \texttt{id} (Number) - Unikalny identyfikator wizyty. \\ \hline
\textbf{Request Body} & 
\vspace{-4mm}
\begin{verbatim}
{
  "visitDate": "2026-05-12T10:30",
  "status": "COMPLETED",
  "room": "Gabinet 112"
}
\end{verbatim}
\vspace{-4mm}
\\ \hline
\textbf{Odpowiedź (Sukces)} & \texttt{200 OK} - \texttt{\{"message": "Wizyta zaktualizowana pomyślnie"\}} \\ \hline
\textbf{Możliwe błędy} & \texttt{404 Not Found} (Gdy wizyta nie istnieje), \texttt{500 Internal Server Error} (Błąd constraintów SQL) \\ \hline
\end{tabularx}

\vspace{1cm}

\subsection{Usuwanie Lekarza}
Wykonuje trwałe usunięcie zasobu lekarza z bazy danych wraz z \textbf{kaskadowym usunięciem} powiązanych wizyt. \index{Kaskadowe Usuwanie}\\

\noindent
\begin{tabularx}{\textwidth}{|l|X|}
\hline
\textbf{Metoda HTTP} & \texttt{DELETE} \\ \hline
\textbf{Ścieżka} & \texttt{/api/doctors/:id} \\ \hline
\textbf{Wymagane role} & \texttt{0} (Admin) \\ \hline
\textbf{Parametry ścieżki} & \texttt{id} (Number) - Unikalny identyfikator lekarza. \\ \hline
\textbf{Logika Biznesowa} & Silnik najpierw egzekwuje \texttt{DELETE FROM visits WHERE doctorId = ?}, zwalniając blokadę \texttt{Foreign Key}, po czym usuwa rekord z tabeli \texttt{doctors}. Następuje zapis encji usuniętej do audytu z odpowiednią flagą. \\ \hline
\textbf{Odpowiedź (Sukces)} & \texttt{200 OK} - \texttt{\{"message": "Lekarz i jego wizyty zostali trwale usunięci z systemu."\}} \\ \hline
\end{tabularx}

\section{Podsumowanie modułu}
Architektura interfejsu programistycznego rygorystycznie chroni integralność relacyjną danych. Middleware autoryzacyjne działają jak strażnik dostępu, a serwis logujący gwarantuje odporność na uszkodzenia operacyjne w sytuacji, gdy rekord źródłowy zostaje usunięty.

% ---------------------------------------------------------
% BIBLIOGRAFIA (Przykładowa)
% ---------------------------------------------------------
\begin{thebibliography}{9}
\bibitem{angular}
  Zespół Google, \emph{Angular Documentation}, 2026.
\bibitem{sqlite}
  D. Richard Hipp, \emph{SQLite Database Architecture}, 2026.
\bibitem{faker}
  Faker.js Community, \emph{FakerPL - Generowanie zbiorów testowych}, 2026.
\end{thebibliography}

% ---------------------------------------------------------
% INDEKS
% ---------------------------------------------------------
\printindex

\end{document}